\section*{第 \ref{ch:b2:viscous-flows} 章——黏性流动}

\begin{solution}{exo:b2:viscous-flows:1}
$\tau = \eta U/h = 0.25 \times 2/2 \times 10^{-4} = \qty{2.5}{kPa}$; $F = \qty{1.25}{kN}$;
$P = Fv = \qty{2.5}{kW}$。干摩擦：$0.3 \times 98 = \qty{29}{N}$——此处
薄而快速剪切的油膜阻力大得多；黏性阻力按
$U/h$ 增长，因此润滑在较低速度、较厚油膜下更为有效（且
它永远不会卡死）。
\end{solution}

\begin{solution}{exo:b2:viscous-flows:2}
(a) $2 \times 10^{-6} \times 3 \times 10^{-5}/10^{-6} = 6 \times 10^{-5}$，蠕动流。(b)
$1.8 \times 1.5/10^{-6} = 2.7 \times 10^6$，湍流。(c) $4 \times 30/1.5 \times 10^{-5}
= 8 \times 10^6$，湍流。(d) $0.025 \times 0.33/3.3 \times 10^{-6} = 2500$，
临界层流。(e) $900 \times 0.5 \times 0.01/0.1 = 45$，层流。
\end{solution}

\begin{solution}{exo:b2:viscous-flows:3}
$D_V = \pi R^4\Delta P/8\eta L = \pi(2 \times 10^{-4})^4 \times 2 \times 10^4/(8 \times 10^{-3}
\times 0.03) = \qty{0.42}{mL/s}$；\qty{12}{s} 可收集 \qty{5}{mL}。$v_{\max} = \Delta PR^2/
4\eta L = \qty{6.7}{m/s}$, $\langle v\rangle = \qty{3.3}{m/s}$, $\mathrm{Re} = 2R\langle v
\rangle/\nu = 1300$：层流。$\tau_w = \Delta PR/2L = \qty{67}{Pa}$。在
\qty{0.8}{mm} 处流量增至十六倍，即 \qty{6.7}{mL/s}——但 $\mathrm{Re} \approx 10^4$：流动
转为湍流，增益较小。
\end{solution}

\begin{solution}{exo:b2:viscous-flows:4}
$v_\infty = 2R^2\Delta\rho\,g/9\eta$。(a) \qty{1.3}{cm/s}，$\mathrm{Re} = 0.02$，
每米 \qty{75}{s}。(b) \qty{0.24}{mm/s}，$\mathrm{Re} = 3 \times 10^{-5}$，\qty{70}{min} 每米。
(c) \qty{1.2}{m/s}，但 $\mathrm{Re} = 16$：斯托克斯公式高估了（实际
速率约 \qty{0.7}{m/s}），每米约一秒。
\end{solution}

\begin{solution}{exo:b2:viscous-flows:5}
(a) $\langle v\rangle = 0.1/0.0707 = \qty{1.4}{m/s}$, $\mathrm{Re} = 4.2 \times 10^5$. (b)
$8\eta LD_V/\pi R^4 = \qty{500}{Pa}$. (c) $\lambda = 0.3/25.5 = 0.0118$, $\Delta P = 0.0118
\times 3333 \times 1000 = \qty{39}{kPa}$：是原来的 $80$ 倍；$P = \Delta PD_V = \qty{3.9}{kW}$。
(d) 速率相同，$\mathrm{Re} = 3 \times 10^5$，$\lambda = 0.0128$，$\Delta P = 0.0128 \times 4717
\times 1000 = \qty{60}{kPa}$，\qty{6}{kW}：单位流量下管壁更多，损失更大。
\end{solution}

\begin{solution}{exo:b2:viscous-flows:6}
(a) $v = 2 \times 25 \times 10^{-12} \times 50 \times 9.81/(9 \times 10^{-3}) = \qty{2.7}{\micro m/s}$;
沉降 \qty{5}{cm} 需 $1.8 \times 10^4$ s，即五小时。(b) \qty{2.7}{mm/s}，\qty{18}{s}。(c)
$\Omega = \sqrt{1000g/r} = \qty{313}{rad/s} = \qty{3000}{rpm}$。(d) $v \propto R^2$，
大颗粒先沉降；依次提高转速离心，
即可按粒径分离。
\end{solution}

\begin{solution}{exo:b2:viscous-flows:7}
(a) $F = 0.6 \times 144 \times 0.35 = \qty{30}{N}$, $+\qty{4}{N}$: $P = 34 \times 12 =
\qty{410}{W}$。(b) $47 + 4 = \qty{51}{N}$，\qty{770}{W}。(c) 相对速率
\qty{17}{m/s}：$61 + 4 = \qty{65}{N}$，乘以地面速率：\qty{780}{W}。(d)
阻力 $\qty{21}{N}$：\qty{300}{W}，节省 \qty{110}{W}。(e) $mg\sin\alpha = \qty{39}{N} =
4 + 0.21v^2$: $v = \qty{13}{m/s}$, \qty{47}{km/h}.
\end{solution}

\begin{solution}{exo:b2:viscous-flows:8}
(a) 高尔夫球：$\tfrac12\rho C_xS = \qty{2.2e-4}{kg/m}$，$v_\infty = \sqrt{0.45/2.2 \times 10^{-4}}
= \qty{45}{m/s}$；跳伞者 $\sqrt{785/0.42} = \qty{43}{m/s}$；降落伞 $\sqrt{785/15} =
\qty{7.2}{m/s}$。(b) $m\dot v = mg - kv^2$，即 $\dot v = g(1 - v^2/v_\infty^2)$，
其从静止出发的解为 $v_\infty\tanh(gt/v_\infty)$。(c) 当 $\tanh x = 0.9$ 时
$x = 1.47$: $t = 1.47 \times 43/9.81 = \qty{6.5}{s}$; $z = (v_\infty^2/g)\ln\cosh x =
190 \times 0.83 = \qty{160}{m}$.
\end{solution}

\begin{solution}{exo:b2:viscous-flows:9}
(a) 剪切速率 $\Omega R/e = 0.5/10^{-3} = \qty{500}{s^{-1}}$；$\Gamma = \tau\cdot2\pi Rh
\cdot R$: $\tau = 0.02/(2\pi \times 2.5 \times 10^{-3} \times 0.1) = \qty{12.7}{Pa}$. (b) $\eta
= 12.7/500 = \qty{2.5e-2}{Pa.s}$. (c) $\mathrm{Re} = \rho Ue/\eta \approx 1000 \times 0.5 \times
10^{-3}/0.025 = 20$：层流。(d) 薄间隙使剪切速率均匀
（平面近似），并使流动保持层流。
\end{solution}

\begin{solution}{exo:b2:viscous-flows:10}
$D_V = \qty{8.3e-5}{m^3/s}$. (a) $R_h = 8\eta L/\pi R^4 = 8 \times 3.5 \times 10^{-3} \times
0.4/(\pi \times 2.4 \times 10^{-8}) = \qty{1.5e5}{Pa.s/m^3}$, $\Delta P = \qty{12}{Pa}$. (b)
单根毛细血管 $\qty{3.5e16}{Pa.s/m^3}$；$10^{10}$ 根并联：\qty{3.5e6}{Pa.s/m^3}；
$\Delta P = \qty{290}{Pa}$（\qty{2}{mmHg}）。(c) 单根小动脉 $8 \times 3.5 \times 10^{-3}
\times 2 \times 10^{-3}/(\pi \times 10^{-20}) = \qty{1.8e15}{Pa.s/m^3}$；整组：
\qty{1.8e7}{Pa.s/m^3}；$\Delta P = \qty{1.5}{kPa}$：小动脉主导了
两者的阻力（压强降发生在小血管，
而非大血管）。(d) $P = 13000 \times 8.3 \times 10^{-5} = \qty{1.1}{W}$，约 $1\%$。(e)
$(1/1.1)^4 = 0.68$：约三分之一的小动脉阻力被去除——
血管扩张剂正是通过这种方式降低血压。
\end{solution}

\begin{solution}{exo:b2:viscous-flows:11}
(a) $F = b\int_0^L\tau_w\,\dd x = 0.332\rho U^2b\sqrt{\nu/U}\int_0^L\dd x/\sqrt x =
0.664\rho U^2b\sqrt{\nu L/U} = 0.664\rho U^2bL/\sqrt{\mathrm{Re}_L}$. (b) $\mathrm{Re}_L =
3 \times 10^5$（层流）；$\delta \approx 5\sqrt{\nu L/U} = \qty{9}{mm}$；$F = 0.664 \times 1000
\times 0.09 \times 1/548 = \qty{0.11}{N}$（每面），共 \qty{0.22}{N}。(c) $\mathrm{Re}_L = 10^7$：
\omterm{prop:b2:viscous-flows:bl}{边界层}在最初几厘米之后转为湍流，层流
公式失效（实际阻力为其数倍）。(d) $\mathrm{Re}_x =
10^9$: $\delta \approx 0.37 \times 100/63 = \qty{0.6}{m}$.
\end{solution}

\begin{solution}{exo:b2:viscous-flows:12}
(a) $0 = \rho g + \eta v''(y)$，$v(0) = 0$（无滑移），$v'(h) = 0$（自由
表面处无应力）。(b) 两次积分：$v = (\rho g/2\eta)(2hy - y^2)$；$q =
\int_0^hv\,\dd y = \rho gh^3/3\eta$. (c) $v_s = \rho gh^2/2\eta = 1200 \times 9.81 \times
4 \times 10^{-8}/1 = \qty{0.47}{mm/s}$；$q = \qty{6.3e-8}{m^2/s}$；十
分钟流出 \qty{28}{cm}——这一厚度的牛顿涂料会流失；实际
涂料呈剪切变稀与触变性，
而且过厚时仍会下垂，因为 $v_s \propto h^2$。(d) $1400 \times 9.81 \times 4 \times 10^{-6}/20 = \qty{2.7}{mm/s}$。
\end{solution}

\begin{solution}{pb:b2:viscous-flows:1}
\textbf{1.} $S = \qty{4.9}{cm^2}$, $\langle v\rangle = 83/4.9 = \qty{0.17}{m/s}$; $\mathrm{Re}
= 1060 \times 0.17 \times 0.025/3.5 \times 10^{-3} = 1300$：层流（脉动峰值
可达 $\sim 5000$）。

\textbf{2.} $v_{\max} = \qty{0.34}{m/s}$；$\tau_w = 4\eta\langle v\rangle/R = \qty{0.19}{Pa}$。

\textbf{3.} $\mathrm{Re} = 1060 \times 3 \times 10^{-4} \times 8 \times 10^{-6}/3.5 \times 10^{-3} =
7 \times 10^{-4}$；$\Delta P = 8\eta L\langle v\rangle/R^2 = 8 \times 3.5 \times 10^{-3} \times 10^{-3}
\times 3 \times 10^{-4}/1.6 \times 10^{-11} = \qty{525}{Pa} = \qty{3.9}{mmHg}$。

\textbf{4.} $1/0.3 = \qty{3.3}{s}$：时间足够长，足以让氧气
扩散到组织中的几微米距离（见\cref{ch:b2:particle-diffusion}）。

\textbf{5.} 单根小动脉的 $8\eta L/\pi R^4 = \qty{1.8e15}{Pa.s/m^3}$；$10^8$ 根
并联：\qty{1.8e7}{Pa.s/m^3}；$\Delta P = 1.8 \times 10^7 \times 8.3 \times 10^{-5} =
\qty{1.5}{kPa} = \qty{11}{mmHg}$。

\textbf{6.} $P = \Delta P\,D_V = 13000 \times 8.3 \times 10^{-5} = \qty{1.1}{W}$，约占
人体功率的 $1\%$。

\textbf{7.} $(1/0.6)^4 = 7.7$。静止时，下游血管扩张并
相互抵消；用力时，需求增至四到五倍，变窄的动脉无法满足，
于是出现心绞痛。

\textbf{8.} 速率变为五倍：$\mathrm{Re} \approx 6500$，湍流——
运动时听诊器能听到的杂音。

\textbf{9.} 在小静脉和毛细血管中缓慢、低切变的流动里，
红细胞聚集，表观\omterm{def:b2:viscous-flows:viscosity}{黏度}升高（主动脉中速率快，情况相反）；
所有对小血管压降的估算
都是下界。

\textbf{10.} $\langle v\rangle = 3/1.13 = \qty{2.65}{m/s}$；$\mathrm{Re} = 850 \times 2.65 \times
1.2/0.02 = 1.35 \times 10^5$：湍流。

\textbf{11.} $\lambda = 0.316/19.2 = 0.0165$；$\Delta P/L = (0.0165/1.2) \times \tfrac12 \times
850 \times 7.0 = \qty{41}{Pa/m} = \qty{41}{kPa/km}$；全线压降为 \qty{534}{bar}。

\textbf{12.} $534/80 = 6.7$：至少需要七个泵站。

\textbf{13.} $P = \Delta P\,D_V = 5.34 \times 10^7 \times 3 = \qty{160}{MW}$；
\qty{53}{MJ/m^3}，占石油能量的 $0.15\%$。

\textbf{14.} $\Delta P = 32\eta L\langle v\rangle/D^2 = 32 \times 0.02 \times 1.3 \times 10^6
\times 2.65/1.44 = \qty{15}{bar}$：少了三十五倍。湍流
代价高昂；百万分之几的长链聚合物即可使其延迟出现，
并使压降降低数十个百分点。

\textbf{15.} $\mathrm{Re} = 1350$：层流；泊肃叶定律：$\Delta P = 32 \times 2 \times 1.3 \times
10^6 \times 2.65/1.44 = \qty{1500}{bar}$——不可能；油必须保持
温暖，因此需要保温和加热。

\textbf{16.} $1.3 \times 10^6/2.65 = \qty{4.9e5}{s} = \qty{5.7}{\text{天}}$。

\textbf{17.} $\tau_w = (\Delta P/L)D/4 = 41 \times 0.3 = \qty{12}{Pa}$；$F = \tau_w\pi DL =
\qty{6e7}{N}$。

\textbf{18.} $\Delta P \propto \lambda\langle v\rangle^2/D \propto D^{-4.75}$：$1.2^{4.75} = 2.4$。

\textbf{19.} $v = 2R^2\Delta\rho\,g/9\eta = 2 \times 2.5 \times 10^{-9} \times 1800 \times 9.81/
0.18 = \qty{0.49}{mm/s}$；$\mathrm{Re} = 2 \times 10^{-3}$；穿过
\qty{1.2}{m} 需要 \qty{41}{min}。油流动时，每秒几厘米的湍流脉动
使其保持悬浮；在停止的管线中，它
不到一小时便沉降下来。

\textbf{20.} $v = 2 \times 2.5 \times 10^{-11} \times 1000 \times 9.81/1.62 \times 10^{-4} =
\qty{3.0}{mm/s}$；$\mathrm{Re} = 2 \times 10^{-3}$；\qty{1}{km} 所需时间为 $3.3 \times 10^5$ s，约
四天：超过每秒几毫米的任何上升气流都能使其悬浮。

\textbf{21.} \qty{1.2}{m/s}，$\mathrm{Re} = 16$：斯托克斯定律大约高估了
两倍左右（实际为 \qty{0.7}{m/s}）。

\textbf{22.} 斯托克斯定律将给出 \qty{120}{m/s} 和 $\mathrm{Re} \sim 10^4$：荒谬。
二次阻力：$\tfrac43\pi R^3\rho_wg = \tfrac12\rho v^2\pi R^2C_x$，$v = \sqrt{8R
\rho_wg/3\rho C_x} = \qty{6.6}{m/s}$（$\mathrm{Re} = 900$，自洽）；从
\qty{1}{km} 处下落需要 \qty{150}{s}。

\textbf{23.} \qty{1}{cm/s} $>$ \qty{3}{mm/s}：液滴上升；毛毛雨滴
需要约 \qty{1}{m/s}，对流云可以提供。

\textbf{24.} $2\gamma/R$：\qty{29}{kPa}、\qty{1.4}{kPa}、\qty{144}{Pa}；\omterm{thm:b2:euler-bernoulli:bernoulli}{动压}
$\tfrac12\rho v^2$ 分别为 $5 \times 10^{-6}$、$0.3$ 和 \qty{26}{Pa}。该比值
在 \qty{2}{mm} 处达到 $0.2$；在 \qty{5}{mm} 处（$v \approx \qty{9}{m/s}$，$\tfrac12\rho v^2 =
\qty{49}{Pa}$ 对比 \qty{58}{Pa}），气\omterm{thm:b2:euler-bernoulli:bernoulli}{动压}力使
液滴压扁并使其破碎：没有雨滴的直径能超过约 \qty{5}{mm}。

\textbf{25.} 主动脉 $10^3$（层流，接近转捩）；毛细血管 $10^{-3}$
（泊肃叶）；小动脉 $10^{-2}$（泊肃叶）；温热管线 $10^5$
（湍流，$\lambda(\mathrm{Re})$）；冷管线 $10^3$（泊肃叶）；油中的砂粒
$10^{-3}$（斯托克斯）；云滴 $10^{-3}$（斯托克斯）；雨滴 $10^3$
（二次阻力，$C_x \approx 0.5$）。
\end{solution}
